\subsection{Deployment, Monitoring, and Feedback Loops}

Deployment turns a fitted model into part of a larger stochastic system.

\begin{definition} \label{def:sl-monitoring} \hlt{Monitoring}\\
Monitoring tracks input distributions, prediction distributions, losses when labels arrive, calibration, latency, missingness, and operational failures.  It detects when the assumptions behind validation no longer hold.
\end{definition}

\begin{definition} \label{def:sl-feedback-loop} \hlt{Feedback Loop}\\
A feedback loop occurs when model predictions influence the future data observed by the model.  Recommendation, credit, policing, and medical triage systems can all change their own training distribution.
\end{definition}

\begin{method} \label{method:sl-model-update-policy} \hlt{Model Update Policy}\\
An update policy specifies when to retrain, what data are eligible, how to validate a replacement, and how to roll back.  It is part of the statistical procedure because it determines the data seen by future models.
\end{method}

\begin{remark} \label{rem:sl-human-system-boundary} \hlt{Human-System Boundary}\\
When a learning model supports human decisions, the deployed predictor is only one component.  The effective policy includes interfaces, incentives, overrides, monitoring, and accountability.
\end{remark}
